
\newpage\selectlanguage{italian}
\chapter*{Ringraziamenti}
\begin{onehalfspacing}
Desidero esprimere la mia più sincera gratitudine a tutte le persone — familiari e amici — che mi hanno accompagnato durante questi ultimi due anni.  
Sono stati anni intensi e ricchi di cambiamenti: oltre all’aspetto universitario che fa da cornice, ci sono state esperienze — come quella di Siemens a Milano, ma anche quelle di supplenza a Benevento e Sant’Agata de’ Goti — che mi hanno profondamente arricchito e aiutato a crescere sia come persona che come futuro professionista.
Ringrazio me stesso per la testardaggine con cui ho portato avanti alcune scelte, anche quando non era facile.  
Ringrazio me stesso per aver avuto il coraggio di lanciarmi in tante cose diverse, a volte forse troppe, chiedendomi spesso se tutto ciò avesse davvero senso.  
Tra queste esperienze, anche la musica ha avuto un ruolo importante: nonostante la costanza ballerina, sono riuscito a suonare la chitarra e il basso, a volte insieme a mio padre, altre volte con amici e persone conosciute lungo il cammino. Sono momenti che hanno dato ritmo e leggerezza a questo percorso.
Credo che il nostro compito, al di là del ruolo che assumiamo — ingegneri, medici, insegnanti o altro — sia prima di tutto quello di essere persone disposte a imparare e a maturare con il tempo, senza fuggire né ignorare i segnali che provengono dagli altri o da noi stessi.  
Siamo una generazione che, nella sua fragilità, trova una grande risorsa: \textbf{il dubbio}. Senza il dubbio non potremmo metterci veramente in gioco.  
Spesso cerchiamo risposte sicure e rassicuranti, che ci leghino a una persona o a un luogo, ma io credo nell’esatto contrario: \textbf{l’incertezza è ciò che ci mantiene vivi}, perché ogni cosa che valga la pena di essere vissuta è inevitabilmente difficile.  
\noindent
Come scriveva Nietzsche:
\begin{displayquote}
«Il deserto cresce: guai a chi cela deserti dentro di sé!»
\end{displayquote}
\\
\noindent
E allora distruggiamolo, questo deserto — con la curiosità, con il coraggio e con la voglia di continuare a imparare.

\end{onehalfspacing}

\vspace{1.5cm}
\selectlanguage{english}
\chapter*{Acknowledgments}
\begin{onehalfspacing}

I would like to express my deepest gratitude to all the people — family and friends — who have accompanied me during these past two years.  
They have been intense and transformative years: beyond the academic aspect that frames them, there have been experiences — such as my time at Siemens in Milan, and my teaching assignments in Benevento and Sant’Agata de’ Goti — that have deeply enriched me and helped me grow both personally and professionally.  
I thank myself for being stubborn enough to stand by certain choices, even when it wasn’t easy.  
I thank myself for having the courage to dive into so many different things — sometimes too many — and for often wondering whether it all truly made sense.  
Among these experiences, music has also played an important role: even without great consistency, I managed to play guitar and bass — sometimes with my father, other times with friends and fellow musicians I met along the way. These moments added rhythm and lightness to my journey.
I believe that our task, beyond any role we might have — engineers, doctors, teachers, and so on — is first and foremost to be people willing to learn and mature over time, without running away or ignoring the signals that come from others or from within ourselves.  
We are a generation that turns its weakness into its greatest resource: \textbf{doubt}.  
Without doubt, we cannot truly challenge ourselves.  We often seek safe and comforting answers that tie us to a person or a place, but I believe the opposite is true: \textbf{uncertainty is all we really have}, and everything worth living for is inevitably difficult.  
\noindent
As Nietzsche once wrote:
\begin{displayquote}
“The desert grows: woe to him who hides deserts within himself.”
\end{displayquote}
\\
\noindent
So let us destroy that desert — with curiosity, with courage, and with the will to keep learning.

\end{onehalfspacing}
